\section{A quantitative answer to Suzuki's density question}
\label{sec:suzuki-density}
For $x>1$ let $E(x)=\sum_{n\le x}'\Lambda(n)/\sqrt n-2\sqrt x$.
The prime on the sum gives half weight to its last term when $x$ is a prime power.
Suzuki explicitly asks for the logarithmic density of $E(x)\le0$ in
\S1.1 after (6) of \emph{On variants of Chebyshev's conjecture}
(\href{https://arxiv.org/html/2411.07436v3}{corrected text}).
His Theorem~1 proves a sign statement for an integral of $E$.
The incoming report \texttt{grh-transfers.zip}, file \texttt{races.tex},
compares primes in distinct residue classes with an additional logarithmic smoothing.
Here we bound how often the \emph{unsmoothed} error itself is negative.

\paragraph{Theorem.}
Put $p(a,L)=L^{-1}|\{t\in[a,a+L]:E(e^t)>0\}|$, where $|\cdot|$
denotes Lebesgue measure. Using RH,
\[
 0<\liminf_{L\to\infty}\inf_{a\ge1}p(a,L)
 \ \le\ \limsup_{L\to\infty}\sup_{a\ge1}p(a,L)<0.00323 .
\]
In particular the lower logarithmic density of $E(x)<0$ exceeds
$99.677\%$, and its upper logarithmic density is strictly below one.
The placement of the logarithmic interval is arbitrary. Global simplicity
and linear independence of zero ordinates are both unnecessary.
The ordinary logarithmic density at this particular threshold is still undetermined by this argument.

\paragraph{The frequencies and their mean square.}
Write $\gamma$ for a distinct positive ordinate of a nontrivial zeta zero and
$m_\gamma$ for its multiplicity. Suzuki's equations (36) and (47) give, under RH,
\begin{align*}
 E(e^t)&=-c+F(t)+r(t),&
 c&=\tfrac12(\pi/2+\log(8\pi)+\gamma_{\rm E}),\\
 F(t)&=-2\lim_{T\to\infty}\sum_{0<\gamma\le T}
                 m_\gamma\frac{\sin(\gamma t)}{\gamma},&
 r(t)&=\sum_{j\ge1}\frac{e^{-(2j+1/2)t}}{2j+1/2}.
\end{align*}
Here $\gamma_{\rm E}$ is Euler's constant and $r(t)$ tends exponentially to zero.
This uses the same weighted-prime-sum/zero-residue correspondence as the
repository's \texttt{New/Reduction/Zeros/PrimeError}, with Suzuki's
weight in place of its one-sided remainder. The geometric determinant
inequality enters this particular deduction through its final RH theorem.

We give the convergence argument because pointwise explicit formulas alone
do not justify the density estimate. Define the Stepanov norm
$\|f\|_{S^2}=\sup_b(\int_b^{b+1}|f(t)|^2dt)^{1/2}$ on the real line.
For a finite exponential sum $P(t)=\sum A_\lambda e^{i\lambda t}$, put
$B_k=\sum_{k\le\lambda<k+1}|A_\lambda|$. Then
\[
                \|P\|_{S^2}^2\ \ll\ \sum_{k\in\mathbb Z} B_k^2 .
\]
Indeed, $W(u)=(\sin(u/2)/(u/2))^2$ dominates a positive constant on
$[-1/2,1/2]$ and its Fourier transform with angular frequencies is bounded
and supported on $[-1,1]$. Integrate $|P|^2$ against each translate of $W$.
Only equal or adjacent unit frequency blocks contribute, and
$2B_kB_{k+1}\le B_k^2+B_{k+1}^2$.
The Riemann--von Mangoldt estimate gives $O(\log(2+k))$ zeros with
multiplicity per unit block at positive height $k$.
For tails of the displayed sine series this gives
$B_k\ll\log(2+|k|)/(1+|k|)$ away from the origin.
The squared bound is summable. The truncations therefore have an $S^2$
limit on the real line, agreeing almost everywhere for $t>0$ with the
explicit formula by local $L^2$ convergence and pointwise convergence.
This is a direct instance of the Stepanov approach of
Kaczorowski and Ramar\'e, \emph{Acta Arith.} 106 (2003), 277--297
(\href{https://doi.org/10.4064/aa106-3-6}{paper}).

For a finite truncation $P$, the means of $P^2$ and $(P_+)^2$ over
intervals of length $L$ converge uniformly in their starting point.
In fact the phase flow is uniformly distributed in its compact orbit closure,
with Haar measure there. This follows on characters by integration of
$e^{i\lambda t}$, and on continuous functions by uniform approximation.
Conjugation on the phase group preserves Haar measure and sends the sine
sum to its negative. Thus the second mean is half the first, including
when ordinates satisfy rational relations.
The positive-part map is $1$-Lipschitz. On any interval of length $L\ge1$,
the normalized $L^2$ norm of the difference between $P$ and $F$ is at most
$\sqrt{1+1/L}\,\|P-F\|_{S^2}$. Taking the truncation limit and using
orthogonality of distinct frequencies yields the uniform mean
\[
          \lim_{L\to\infty}\frac1L\int_a^{a+L} (F(t)_+)^2dt
          =\frac V2,\qquad V=2\sum_{\gamma>0}\frac{m_\gamma^2}{\gamma^2}.
\]
For any $0<\eta<c$, the inequality $E(e^t)>0$ forces
$F(t)>c-\eta$ after a fixed starting time. The initial segment contributes
$O_\eta(1/L)$ uniformly for $a\ge1$.
Markov's inequality, followed by $\eta\downarrow0$, gives $V/(2c^2)$
as an upper bound for the upper density.

\paragraph{An explicit multiplicity allowance.}
Put $H=3\cdot10^{12}$ and $S=2+\gamma_{\rm E}-\log(4\pi)$.
The Hadamard product for $\xi$, on subtracting its logarithmic derivatives at
$1$ and $0$, gives
\[
 S=2\xi'(1)/\xi(1)=2\sum_{\gamma>0}\frac{m_\gamma}{1/4+\gamma^2}.
\]
Here the first evaluation follows from the Laurent expansion of $\zeta$ at
one and the gamma factor. For $u=(1/4+\gamma^2)^{-1}$ one has $u\le S/2$
and $\gamma^{-2}=u/(1-u/4)$. Thus
$2\sum m_\gamma/\gamma^2\le S/(1-S/8)$.
Platt and Trudgian verified simplicity through $H$
(\href{https://doi.org/10.1112/blms.12460}{Bull. London Math. Soc. 53 (2021)}).
At larger heights the zero-count estimate of Hasanalizade, Shen and Wong
(\href{https://arxiv.org/abs/2107.06506}{Corollary 1.2}) gives
\[
 \left|N(t)-\frac{t}{2\pi}\log\frac{t}{2\pi e}\right|
 \le 0.1038\log t+0.2573\log\log t+9.3675 <\log t \quad(t\ge H).
\]
Consequently $N(t)\le t\log t$ and its jump at $\gamma>H$ is at most
$2\log\gamma$. The extra multiplicities in $V$ contribute at most
\[
 4\int_{(H,\infty)}\frac{\log t}{t^2}\,dN(t)
 \le 8\int_H^\infty \frac{\log^2t}{t^2}\,dt
 =\frac{8}{H}\big((\log H)^2+2\log H+2\big).
\]
Integration by parts drops a nonpositive lower boundary term.
It follows that $V<0.04647$, $c>2.686$, and $V/(2c^2)<0.00323$.

\paragraph{Why the other sign has positive lower density.}
The classical Littlewood oscillation theorem makes
$e^{-t/2}(\psi(e^t)-e^t)$ unbounded above and below. Its zero series
has coefficients $-1/(1/2+i\gamma)$, differing from those of $F$ by
$O(\gamma^{-2})$ per zero. The difference series is uniformly bounded,
as are the remaining terms at $t\ge1$. Thus $E(e^t)$ also takes arbitrarily
large positive and negative values (as Suzuki observes in \S1.1).
The difference between the right-continuous and the half-weight convention
tends to zero. Between consecutive prime powers $E(x)$ is smooth, and at
a jump its one-sided values bracket the half-weight value.
Choose a sufficiently late interval $I$ of length at most one and
$\eta>0$ such that $F>c+2\eta$ almost everywhere on $I$.

The $S^2$ limit of finite trigonometric polynomials has relatively dense
$S^2$ almost periods: approximate it by one polynomial, then use compactness
of that polynomial's phase group to obtain relatively dense uniform almost
periods. For every such $\tau$ with
$\|F(\cdot+\tau)-F\|_{S^2}<\eta\sqrt{|I|/2}$, at least half of $I+\tau$
has $F>c+\eta$, by Chebyshev's inequality. These translates occur with
bounded gaps. More explicitly, if every shift interval of length $D$ contains
such a $\tau$, each consecutive block of length $D+|I|$ contains one
complete translated interval. Since $r(t)>0$, this proves the stated
strictly positive uniform lower density. Finally $E(x)=0$ on a countable
set because it is strictly decreasing between prime powers. This permits
either sign convention at the density threshold.
